\documentclass[11pt]{article}
\usepackage[margin=0.82in]{geometry}
\usepackage{amsmath,amssymb,amsthm,microtype}
\usepackage[hidelinks]{hyperref}

\title{Separable Hilbert spaces need not have the nice-geodesics structure}
\author{Full counterexample for arXiv:2204.01461}
\date{13 August 2026}

\newtheorem{theorem}{Theorem}
\newtheorem{corollary}{Corollary}
\newcommand{\N}{\mathbb N}

\begin{document}
\maketitle

\begin{center}
\fbox{\parbox{0.91\textwidth}{\textbf{Status: full negative answer.}
Every infinite-dimensional separable Hilbert space fails the source's
nice-geodesics condition, for every choice of countable dense test set.}}
\end{center}

\section*{Source question}

Let $H$ be a separable Hadamard space and let $(y_n)_{n\in\N}$ be dense in
$H$.  The paper calls the geodesic structure \emph{nice} if, for every
$x\in H$ and every sequence $(x_k)$,
\[
 P_{[x,y_n]}x_k\longrightarrow x\quad\text{for every }n
 \quad\Longrightarrow\quad
 P_\gamma x_k\longrightarrow x\quad\text{for every geodesic segment }
 \gamma\text{ starting at }x.
\]
After using this as a hypothesis in its compactness theorem, the paper asks:
``Does every separable Hadamard space enjoy the nice geodesics structure?''

\section*{Counterexample theorem}

\begin{theorem}\label{thm:counterexample}
Let $H$ be an infinite-dimensional separable real Hilbert space.  For every
countable dense sequence $(y_n)$ in $H$, there exist a unit vector $u$ and an
unbounded sequence $(z_k)$ such that
\[
 P_{[0,y_n]}z_k\longrightarrow0\quad(n\in\N),
 \qquad
 P_{[0,u]}z_k=u\quad(k\in\N).
\]
Consequently $H$, and in particular $\ell^2$, does not have the
nice-geodesics structure.
\end{theorem}

\begin{proof}
Set $E_k=\operatorname{span}\{y_1,\ldots,y_k\}$.  Each $E_k$ is a
finite-dimensional closed nowhere-dense subset of $H$.  By the Baire category
theorem, choose a unit vector
\[
 u\notin\bigcup_{k=1}^\infty E_k.
\]
Let $Q_k$ be the orthogonal projection onto $E_k$, put
$v_k=u-Q_ku\ne0$, and define
\[
 z_k=\frac{v_k}{\|v_k\|^2}.
\]
For $n\le k$, we have $z_k\perp E_k$ and hence
$\langle z_k,y_n\rangle=0$.  The nearest point to $z_k$ on the segment
$[0,y_n]=\{ty_n:0\le t\le1\}$ is therefore exactly $0$.  Thus, for each
fixed $n$, $P_{[0,y_n]}z_k=0$ for every $k\ge n$.

On the other hand, $v_k\perp Q_ku$, so
\[
 \langle z_k,u\rangle
 =\frac{\langle v_k,Q_ku+v_k\rangle}{\|v_k\|^2}=1.
\]
For $0\le t\le1$,
\[
 \|z_k-tu\|^2=\|z_k\|^2-2t+t^2,
\]
whose unique minimum occurs at $t=1$.  Hence
$P_{[0,u]}z_k=u$ for every $k$, contradicting the defining implication.

Finally, density of $(y_n)$ makes $\bigcup_kE_k$ dense in $H$, whence
$Q_ku\to u$.  Thus $\|v_k\|\to0$ and
$\|z_k\|=\|v_k\|^{-1}\to\infty$.
\end{proof}

\section*{Sharp Hilbert-space classification}
\vspace{0.35em}

\begin{theorem}\label{thm:finite}
Every finite-dimensional Hilbert space has the nice-geodesics structure,
for every countable dense test set.
\end{theorem}

\begin{proof}
Translate the base point to $0$.  Suppose
$P_{[0,y_n]}z_k\to0$ for every $n$.  First, $(z_k)$ is bounded.  Otherwise,
after passing to a subsequence, compactness of the unit sphere gives
$z_k/\|z_k\|\to w$ for a unit vector $w$.  Choose a test point $y_n$ close
enough to $w$ that $\langle w,y_n\rangle>0$.  Then
$\langle z_k,y_n\rangle\to+\infty$, so eventually the projection of $z_k$
onto $[0,y_n]$ is the nonzero endpoint $y_n$, a contradiction.

Let $M=\sup_k\|z_k\|$.  For a nonzero $y$, the projection onto $[0,y]$ is
\[
 P_{[0,y]}z=
 \min\!\left\{1,\max\!\left\{0,
 \frac{\langle z,y\rangle}{\|y\|^2}\right\}\right\}y.
\]
On the compact set $\{z:\|z\|\le M\}$ this expression depends uniformly
continuously on $y$ near any fixed nonzero endpoint.  Approximate an arbitrary
$y$ by one fixed test point $y_n$, then let $k\to\infty$.  It follows that
$P_{[0,y]}z_k\to0$.  The zero endpoint is trivial, so the defining implication
holds for every segment.
\end{proof}

\begin{corollary}
Among separable Hilbert spaces, the nice-geodesics property is equivalent to
finite dimensionality.
\end{corollary}

\section*{Scope and priority check}

Theorem~\ref{thm:counterexample} completely answers the source's universal
question in the negative; no additional curvature or nonlinear construction
is needed.  It does not decide the paper's separate question about weak
compactness versus weak sequential compactness for arbitrary unbounded sets.

The run indexes and bounded primary-source searches through 13 August 2026
found no later answer under the paper's terminology.  The closely related
same-day paper arXiv:2204.01291 analyzes weak topologies for unbounded nets and
the nonsequential weak topology of infinite-dimensional Hilbert space, but it
does not state this countable-segment counterexample.  Priority should still
be checked by a specialist before external dissemination.

\end{document}
